\subsection{DISTRIBUTIVITY FORMULAE}

PROPOSITION 8. Let \(((\mathbf{X}_{\lambda,1})_{i\in \mathbf{J}_2})_{\lambda \in \mathbf{L}}\) be a family (with index set L) of families of sets. Suppose that \(\mathbf{L}\neq \emptyset\) and that \(J_{\lambda}\neq \emptyset\) for each \(\lambda \in \mathbf{L}\) . Let

\[
\mathrm{I} = \prod_ {\lambda \in \mathrm{L}} \mathrm{J} _ {\lambda} \neq \emptyset .
\]

Then we have

\[
\begin{array}{l} \bigcup_ {\lambda \in \mathbf {L}} \Big (\bigcap_ {t \in \mathbf {J} _ {\lambda}} \mathbf {X} _ {\lambda , t} \Big) = \bigcap_ {f \in \mathbf {I}} \Big (\bigcup_ {\lambda \in \mathbf {L}} \mathbf {X} _ {\lambda , f (\lambda)} \Big), \\ \bigcap_ {\lambda \in \mathbf {L}} \Big (\bigcap_ {t \in \mathbf {J} _ {\lambda}} \mathbf {X} _ {\lambda , t} \Big) = \bigcup_ {f \in \mathbf {I}} \Big (\bigcap_ {\lambda \in \mathbf {L}} \mathbf {X} _ {\lambda , f (\lambda)} \Big) \end{array}
\]

(``distributivity'' of union over intersection, and of intersection over union).

Let \(x\) be an element of \(\bigcup_{\lambda \in L} \left( \bigcap_{i \in J_{\lambda}} X_{\lambda, i} \right)\) and let \(f\) be any element of 1. There exists an index \(\lambda\) such that \(x \in \bigcap_{i \in J_{\lambda}} X_{\lambda, i}\) ; consequently \(x \in X_{\lambda, f(\lambda)}\) and hence

\[
x \in \bigcup_ {\lambda \in \mathbf {L}} \mathrm{X} _ {\lambda , f (\lambda)}.
\]

Since this is true for each \(f \in I\) , we have

\[
x \in \bigcap_ {f \in I} \left(\bigcup_ {\lambda \in L} X _ {\lambda , f (\lambda)}\right).
\]

Conversely, let \(x\) be an object which does not belong to the set

\[
\bigcup_ {\lambda \in L} \left(\bigcap_ {i \in J _ {\lambda}} X _ {\lambda , i}\right).
\]

Then for each \(\lambda \in L\) we have \(x \notin \bigcap_{\iota \in J_{\lambda}} X_{\lambda, \iota}\) , which means that the set \(J_{\lambda}'\) of indices \(\iota \in J_{\lambda}\) such that \(x \notin X_{\lambda, \iota}\) is not empty for each \(\lambda \in L\) . By Proposition 5, Corollary 2, there exists a functional graph \(f\) with domain \(L\) such that for each \(\lambda \in L\) we have \(f(\lambda) \in J_{\lambda}'\) . Consequently \(f \in I\) and \(x \notin X_{\lambda, f(\lambda)}\) for each \(\lambda \in L\) ; hence

\[
x \notin \bigcup_ {\lambda \in \mathbf {L}} \mathrm{X} _ {\lambda , f (\lambda)}
\]

and thus

\[
x \notin \bigcap_ {f \in I} \left(\bigcup_ {\lambda \in L} X _ {\lambda , f (\lambda)}\right).
\]

This completes the proof of the first formula. The second formula follows by applying the first to the family \(\left((\mathbb{C}_{\Lambda}X_{\lambda,t})_{t\in J_{\lambda}}\right)_{\lambda \in L}\) , where A denotes the union \(\bigcap_{\lambda \in L}\left(\bigcap_{t\in J_{\lambda}}X_{\lambda}\right)\) .

COROLLARY. Let \((\mathbf{X}_i)_{t\in \mathbf{I}}\) and \((\mathbf{Y}_x)_{x\in \mathbb{K}}\) be two families of sets with non-empty index sets I, K. Then

\[
\begin{array}{l} \left(\bigcap_ {i \in I} X _ {i}\right) \cup \left(\bigcap_ {x \in K} Y _ {x}\right) = \bigcap_ {(i, x) \in I \times K} (X _ {i} \cup Y _ {x}), \\ \left(\bigcup_ {i \in I} X _ {i}\right) \cap \left(\bigcup_ {x \in K} Y _ {x}\right) = \bigcup_ {(i, x) \in I \times K} (X _ {i} \cap Y _ {x}). \end{array}
\]

Let \(\alpha, \beta\) be two distinct objects; apply the formulae of Proposition 8 (with \(L = \{\alpha, \beta\}\) , \(J_{\alpha} = I\) , \(J_{\beta} = K\) ) to the family \(((Z_{\lambda, \mu})_{\mu \in J_{\lambda}})_{\lambda \in L}\) , where \(Z_{\alpha, \iota} = X_{\iota}\) for all \(\iota \in I\) and \(Z_{\beta, x} = Y_{x}\) for each \(x \in K\) . By the existence of the canonical bijection of \(\prod_{\lambda \in L} J_{\lambda}\) onto \(I \times K\) (no. 3) and by Proposition 1 of §4 we obtain the formulae stated.

PROPOSITION 9. Let \(((\mathbf{X}_{\lambda, t})_{t \in \mathbf{J}_{\lambda}})_{\lambda \in \mathbf{L}}\) be a family (with index set L) of families of sets. Let \(\mathbf{I} = \prod_{\lambda \in \mathbf{J}_{\lambda}} \mathbf{J}_{\lambda}\) . Then

\[
\prod_ {\lambda \in L} \left(\bigcup_ {i \in J _ {\lambda}} X _ {\lambda , i}\right) = \bigcup_ {f \in I} \left(\prod_ {\lambda \in L} X _ {\lambda , f (\lambda)}\right)
\]

and (if \(\mathbf{L} \neq \phi\) and \(\mathbf{J}_{\lambda} \neq \phi\) for each \(\lambda \in \mathbf{L}\) )

\[
\prod_ {\lambda \in \mathrm{L}} \left(\bigcap_ {i \in \mathrm{J} _ {\lambda}} \mathrm{X} _ {\lambda , i}\right) = \bigcap_ {f \in \mathrm{I}} \left(\prod_ {\lambda \in \mathrm{L}} \mathrm{X} _ {\lambda , f (\lambda)}\right)
\]

(``distributivity'' of product over union and over intersection).

The first formula is trivially true if \(L = \emptyset\) or if \(J_{\lambda} = \emptyset\) for some \(\lambda \in L\) . If not, let \(g\) be an element of \(\prod_{\lambda \in L} \left( \bigcap_{t \in J_{\lambda}} X_{\lambda, t} \right)\) . For each \(\lambda \in L\) there exists an index \(t \in J_{\lambda}\) such that \(g(\lambda) \in X_{\lambda, t}\) ; in other words, the set \(H_{\lambda}\) of indices \(t \in J_{\lambda}\) such that \(g(\lambda) \in X_{\lambda, t}\) is not empty. By Corollary 2 to Proposition 5 there is therefore a functional graph \(f\) with domain \(L\) such that

\[
f (\lambda) \in \mathrm{H} _ {\lambda}
\]

for each \(\lambda \in L\) , i.e., \(g(\lambda) \in X_{\lambda, f(\lambda)}\) . Hence we have \(g \in \prod_{\lambda \in L} X_{\lambda, f(\lambda)}\) and consequently \(g \in \bigcup_{f \in I} \left( \prod_{\lambda \in L} X_{\lambda, f(\lambda)} \right)\) . Conversely, if

\[
g \in \bigcup_ {f \in \mathbf {I}} \left(\prod_ {\lambda \in \mathbf {L}} \mathrm{X} _ {\lambda , f (\lambda)}\right),
\]

there is a functional graph \(f \in I\) such that for every \(\lambda \in L\) we have

\[
g (\lambda) \in \mathbf {X} _ {\lambda , f (\lambda)}
\]

and, a fortiori, \(g(\lambda) \in \bigcup_{t \in J_1} X_{\lambda,t}\) . This completes the proof of the first formula. The proof of the second formula is analogous but simpler, and we leave it to the reader.

COROLLARY 1. Suppose that \(\mathbf{L} \neq \emptyset\) and that \(J_{\lambda} \neq \emptyset\) for each \(\lambda \in L\) . If for \(e\) achindex \(\lambda \in L\) the family \((X_{\lambda, t})_{t \in J_{\lambda}}\) is a partition of \(X_{\lambda} = \bigcup_{t \in J_{\lambda}} X_{\lambda, t}\) , then the family \(\left( \prod_{\lambda \in L} X_{\lambda, f(\lambda)} \right)_{f \in I}\) is a partition of \(\prod_{\lambda \in L} X_{\lambda}\) .

If we set

\[
\mathrm{P} _ {f} = \prod_ {\lambda \in \mathbf {L}} \mathrm{X} _ {\lambda , f (\lambda)},
\]

then, by virtue of the first formula of Proposition 9, it is sufficient to show that \(\mathbf{P}_f \neq \emptyset\) for all \(f \in \mathbf{I}\) and that \(\mathbf{P}_f \cap \mathbf{P}_g = \emptyset\) whenever \(f\) and \(g\) are distinct elements of \(\mathbf{I}\) . The first point follows from Proposition 5, Corollary 2. As to the second, if \(f \neq g\) , there exists \(\lambda \in \mathbf{L}\) such that

\[
f (\lambda) \neq g (\lambda)
\]

and therefore, by virtue of the hypothesis, \(\mathbf{X}_{\lambda,f(\lambda)} \cap \mathbf{X}_{\lambda,g(\lambda)} = \varnothing\) . It follows that there is no graph belonging to \(P_{f} \cap P_{g}\) ; for if G were such a graph, we would have \(\mathbf{G}(\lambda) \in \mathbf{X}_{\lambda,f(\lambda)} \cap \mathbf{X}_{\lambda,g(\lambda)} = \varnothing\) , which is absurd.

COROLLARY 2. Let \((\mathbf{X}_i)_{i\in \mathbf{I}}\) and \((\mathbf{Y}_x)_{x\in \mathbb{K}}\) be two families of sets. Then

\[
\left(\bigcup_ {i \in I} X _ {i}\right) \times \left(\bigcup_ {x \in K} Y _ {x}\right) = \bigcup_ {(i, x) \in I \times K} (X _ {i} \times Y _ {x})
\]

and, if I and K are non-empty,

\[
\left(\bigcap_ {i \in I} X _ {i}\right) \times \left(\bigcap_ {x \in K} Y _ {x}\right) = \bigcap_ {(i, x) \in I \times K} (X _ {i} \times Y _ {x}).
\]

The proof follows the pattern of the proof of the Corollary to Proposition 8.

PROPOSITION 10. Let \((\mathbf{X}_{t,x})_{(t,x)\in \mathbf{I}\times \mathbf{K}}\) be a family of sets whose index set is the product of two sets I and K. If \(\mathbf{K}\neq \emptyset\) , we have

\[
\bigcap_ {x \in \mathbf {K}} \left(\prod_ {i \in I} X _ {i, x}\right) = \prod_ {i \in I} \left(\bigcap_ {\mathbf {K} \in x} X _ {i, x}\right).
\]

Both sides of the equality to be proved are functional graphs. A graph \(f\) belongs to the left-hand side if and only if, for each \(x \in \mathbf{K}\) , \(f \in \prod_{i \in I} X_{i,x}\) ; that is, if and only if \(f(t) \in X_{t,x}\) for all \((t,x) \in I \times K\) . For \(f\) to belong to the right-hand side it is necessary and sufficient that \(f(t) \in \bigcap_{x \in K} X_{t,x}\) for each \(t \in I\) , i.e., that \(f(t) \in X_{t,x}\) for each pair \((t, x) \in I \times K\) . This completes the proof.

COROLLARY. Let \((\mathbf{X}_t)_{t\in \mathbf{I}}\) and \((\mathbf{Y}_t)_{t\in \mathbf{I}}\) be two families of sets with the same index set \(\mathbf{I}\neq \emptyset\) . Then

\[
\begin{array}{l} \left(\prod_ {i \in I} X _ {i}\right) \cap \left(\prod_ {i \in I} Y _ {i}\right) = \prod_ {i \in I} (X _ {i} \cap Y _ {i}), \\ \left(\bigcap_ {i \in I} X _ {i}\right) \times \left(\bigcap_ {i \in I} Y _ {i}\right) = \bigcap_ {i \in I} (X _ {i} \times Y _ {i}). \end{array}
\]

Apply Proposition 10 to the case where K (resp. I) is a set consisting of two distinct elements.

